% =====================================================================
% Strategy reachability / no-go analysis  --  standalone section
% Generated 2026-09-09 from paper/9-8_wxie/20260908_strategy_reachability/
%
% This file is self-contained and does NOT modify root_wxie.tex.
% To use it, copy the figure files from
%     paper/9-8_wxie/20260908_strategy_reachability/figures/
% into the paper's plot/ directory and add
%     % =====================================================================
% Strategy reachability / no-go analysis  --  standalone section
% Generated 2026-09-09 from paper/9-8_wxie/20260908_strategy_reachability/
%
% This file is self-contained and does NOT modify root_wxie.tex.
% To use it, copy the figure files from
%     paper/9-8_wxie/20260908_strategy_reachability/figures/
% into the paper's plot/ directory and add
%     % =====================================================================
% Strategy reachability / no-go analysis  --  standalone section
% Generated 2026-09-09 from paper/9-8_wxie/20260908_strategy_reachability/
%
% This file is self-contained and does NOT modify root_wxie.tex.
% To use it, copy the figure files from
%     paper/9-8_wxie/20260908_strategy_reachability/figures/
% into the paper's plot/ directory and add
%     % =====================================================================
% Strategy reachability / no-go analysis  --  standalone section
% Generated 2026-09-09 from paper/9-8_wxie/20260908_strategy_reachability/
%
% This file is self-contained and does NOT modify root_wxie.tex.
% To use it, copy the figure files from
%     paper/9-8_wxie/20260908_strategy_reachability/figures/
% into the paper's plot/ directory and add
%     \input{20260908-strategy_reachability}
% where Sec. III.C (the current no-go subsection) sits.
%
% It is written to REPLACE the retained proxy-model subsection
% (\label{sec:nogo_analysis}) and its Fig. 7.  Labels carry the suffix
% _v2 so the two can coexist while the swap is made.
%
% Figure 1 below is a drop-in REPLACEMENT for fig:grasp_strategies: it
% keeps img/grasp_strategies_complex.png and adds two no-go strips whose
% three columns line up with the three photographed strategies.
%
% Numbers here are produced by:
%   code/calibrate_capsule_envelope.py --out results/calibration
%   code/run_reachability_sweep.py --sampler paper-conditions \
%       --contact-policy contact-fingers --floor-model mesh \
%       --floor-tolerance-mm 0 --tag paper_conditions_meshfloor
%   code/run_nogo_volume.py --object paper_strawberry --region quarter-ball \
%       --full-grid --floor-model mesh --floor-tolerance-mm 0
%   code/render_strips.sh   (render_nogo_volume.py --strip --panels-dir,
%                            then compose_nogo_figure.py for the lettered grid)
% Calibration and ablations: results/calibration/README.md
% =====================================================================

\subsection{Strategy-Conditioned Approach Feasibility}
\label{sec:nogo_analysis_v2}

A terminal grasp configuration does not specify a collision-free motion to that
configuration. Because the RH56 finger linkages are coupled, a commanded width
also sets fingertip position and the hand pose required to hold a grasp plane, so
the admissible set depends on which closure strategy issues the command. We
model each strategy as an ordered sequence of waypoints in hand-base position,
hand-base rotation, and the six actuator commands, and check the whole commanded
motion---approach and closure---against the object and the support surface.

Let $\mathbf{x}$ be the world hand-base position at which the motion starts and
let $\sigma$ index the strategies. Writing $A(\mathbf{x},\sigma)=1$ when the
commanded motion is admissible, we report
\begin{equation}
  \nu(\sigma) = \frac{1}{|\mathcal{X}|}\sum_{\mathbf{x}\in\mathcal{X}} A(\mathbf{x},\sigma),
  \label{eq:admissible_fraction}
\end{equation}
the fraction of sampled start poses admitted for that strategy. The sample
$\mathcal{X}$ is the ten approach conditions per object used on hardware---one
adversarial pose level with the grasp point, one overhead, and eight at
\SI{150}{\milli\metre} and \SI{250}{\milli\metre} above it with
$\pm\SI{30}{\milli\metre}$ of lateral jitter---over the same 15 objects, giving
150 trials per strategy.

Admissibility requires three geometric conditions. No part of the hand may
penetrate the support plane; no part of the hand outside the planned contact
fingers may penetrate the object; and the antipodal contact midpoint reached at
the end of the motion must lie inside the object. The planned width is the
object width plus \SI{20}{\milli\metre}, so a correctly executed grasp leaves
about \SI{10}{\milli\metre} of clearance per side and any contact in this model
is unintended. The three strategies differ only in what the wrist does during
closure: naive closure holds the wrist at the pose that straddled the object with
the hand open, iterative closure moves the wrist along the analytical width
sequence, and reflex closure holds the wrist at the final analytical pose while
only the opposing fingers move.

The support-plane test is run against the hand's own collision geometry---palm
boxes and cylinder, phalanx meshes---so it carries no clearance tolerance. A
faster 13-capsule proxy of the hand is used for the object tests, and its
vertical over-approximation was measured rather than assumed: over the 41{,}250
hand states the three strategies command, the lowest point of the capsule proxy
lies a median of \SI{3.4}{\milli\metre} below the lowest point of the collision
geometry (5th to 95th percentile \SI{1.8}{}--\SI{6.8}{\milli\metre}), and it lies
below it at every state, so the capsule test is one-sided. The capsule that
reaches lowest is a distal phalanx in 96.6\% of those states. Because MuJoCo
collides a mesh as its convex hull, the same measurement was repeated against the
undecimated visual meshes; the two agree to a median of
\SI{-0.11}{\milli\metre}, so the collision geometry is the part.

\begin{figure}[!htb]
  \centering
  \includegraphics[width=0.95\columnwidth]{img/grasp_strategies_complex.png}\\[3pt]
  % Both objects and all four strategies in one lettered grid; produced by
  % code/render_strips.sh (render_nogo_volume.py --panels-dir, then
  % compose_nogo_figure.py). The three-strategy strips
  % plot/20260908-nogo_strip_{strawberry,can}.png remain available if the
  % GraspGen-X column is dropped.
  \includegraphics[width=0.95\columnwidth]{plot/20260908-nogo_grid_strawberry_can.png}
  \caption{\textbf{Top:} the three grasp strategies---iterative (left), reflex
  (centre), naive (right). In the iterative approach the hand rotates downward
  from a pre-grasp pose to reach the desired grasp pose. In the reflex approach
  the fingers remain open until the thumb makes contact with the object. In the
  naive approach the fingers are treated like a parallel jaw gripper prior to
  closure, leading to ground collisions for small objects.
  \textbf{Bottom:} start poses each strategy can and cannot execute from, one
  column per strategy with GraspGen-X at its best of 10 ranked candidates, for
  (A) the strawberry, grasp point \SI{15}{\milli\metre} above the table, and
  (B) the can, \SI{110}{\milli\metre}. The sampled region is the quarter ball
  of radius \SI{300}{\milli\metre} above the grasp point and on the approach
  side, the poses an arm is commanded from; the red solid is the subset from
  which the strategy cannot execute and the green shell is the rest, so the two
  together are the region each percentage is a fraction of. The object is drawn
  opaque with a post through it. At the strawberry, clearance during closure
  decides the outcome and the strategies separate: reflex closure leaves 53.3\%
  of the region unreachable, iterative closure 76.9\%, and naive closure and
  GraspGen-X all of it. At the can the support surface is never the binding
  constraint and the three hand-designed sets nearly coincide, 34.0\%, 32.1\%
  and 34.0\%, while GraspGen-X leaves 67.7\% unreachable at every sampled
  height, a restriction on approach direction rather than on clearance.
  Searching all 100 proposals leaves the strawberry region unreachable, since
  every proposal places the hand in the table at the proposed pose itself, and
  about 3\% of the can region.
  Surfaces are the half level of the \SI{30}{\milli\metre} sweep grid, refined
  and low-passed, so the smoothness is presentational and the boundary is
  resolved to \SI{30}{\milli\metre}.}
  \label{fig:grasp_strategies_v2}
\end{figure}

Naive closure is admissible in 48.0\% of trials, iterative closure in 71.3\% and
reflex closure in 90.7\%, against measured hardware success rates of 48.0\%,
82.0\% and 86.7\% over 150 trials each. Naive and reflex closure fall inside the
hardware 95\% confidence intervals (39.8--56.3\% and 80.2--91.7\%); iterative
closure is 10.7 points low. No parameter in the model is fitted to a success
rate.

The naive failures separate into two mechanisms, and neither depends on the start
pose. Terminal misalignment accounts for 60 of 150 trials and occurs for no other
hand-designed strategy: holding the wrist fixed through closure leaves the
coupled fingertip travel uncompensated, displacing the contact midpoint by a
median \SI{17.4}{\milli\metre} (\SI{16.3}{}--\SI{18.6}{\milli\metre}), which
carries it outside every object of grasp width \SI{25}{\milli\metre} or less.
Iterative and reflex closure reach their targets to within
\SI{0.004}{\milli\metre} by construction. The result is a step in grasp-point
height rather than a gradient: naive closure admits none of the 70 trials on
objects whose grasp point lies at or below \SI{15}{\milli\metre} and 72 of the 80
at or above \SI{29}{\milli\metre}, and its rate is identical at both elevated
approach heights (53.3\% at \SI{150}{\milli\metre} and at
\SI{250}{\milli\metre}) because the mechanism is set by object geometry.
Support-surface collision is recorded in 76 naive trials, at a median
penetration of \SI{7.9}{\milli\metre}. The two failure modes overlap rather
than add: every one of the 60 misaligned trials also collides with the support
surface, the index intermediate phalanx reaching the plane in 54 of them at a
median depth of \SI{6.1}{\milli\metre}, so the 78 unreachable naive trials are
the 76 support collisions plus 2 premature object contacts with no collision.
The counts in Table~\ref{tab:ggx_admissibility} are of trials showing each mode
and not a partition of the failures.

Reflex closure retains the low-clearance region that iterative closure loses, and
the advantage is confined to the low approach. At the \SI{250}{\milli\metre}
approach height the two are identical at 100.0\%; at \SI{150}{\milli\metre}
reflex closure holds 100.0\% while iterative closure falls to 53.3\%, the same
rate as naive. Reflex closure is also the only hand-designed strategy admitting
any trial at the adversarial condition level with the grasp point (1 of 15,
against 0 and 0). Its 12 support-plane collisions are uniformly deep---the
shallowest is \SI{51.5}{\milli\metre} and the median
\SI{100.9}{\milli\metre}---and correspond to poses no clearance allowance
recovers, whereas 13 of iterative closure's 41 collisions are grazes no deeper
than \SI{1.7}{\milli\metre}.

That graze population is where the 10.7-point iterative residual sits. Forgiving
support-plane contact shallower than \SI{2}{\milli\metre} moves iterative closure
to 79.3\% and shallower than \SI{5}{\milli\metre} to 84.7\%, bracketing the
measured 82.0\%, on the argument that a silicone fingertip deflects at that depth.
Naive and reflex closure do not move, because neither has a shallow collision to
forgive. Testing the argument needs the per-trial failure cause from the hardware
benchmark, which the archive does not record, so the primary configuration
forgives nothing.

\subsection{A Learned Generator Under the Same Test}
\label{sec:ggx_reachability_v2}

The same model accepts a learned grasp generator without modification.
GraspGen-X proposes a 6-DoF hand pose rather than an aperture, and its execution
has the shape of the wrist-fixed strategies: approach the proposed pose with the
hand open, then interpolate the descriptor's closure with the wrist held still.
We run the released model with the official Inspire-hand descriptor over the same
15 object proxies, producing 100 ranked candidates each, and test candidates in
descending confidence order; a start pose is admissible when any of the 100
executes cleanly, and a pose with no clean candidate is a planning failure, which
is the hardware protocol of Sec.~\ref{sec:vision_evaluation}. Its
terminal check cannot be the antipodal-midpoint test, which presumes a precision
pinch, so we use the descriptor's own declared capture volume and ask whether the
proposal engages the object on its terms.

\begin{table}[t]
  \centering
  \caption{All four strategies over 15 objects and ten approach conditions,
  support plane checked against the hand's collision geometry with no clearance
  tolerance. A trial can trigger more than one failure mode.}
  \label{tab:ggx_admissibility}
  \setlength{\tabcolsep}{4pt}
  \begin{tabular}{lcccc}
    \toprule
    & & Support & Premature & Terminal \\
    Strategy & Admissible & collision & contact & miss \\
    \midrule
    Naive       & 72/150 (48.0\%)  & 76 & 10 & 60 \\
    Iterative   & 107/150 (71.3\%) & 41 & 10 &  0 \\
    Reflex      & 136/150 (90.7\%) & 12 &  5 &  0 \\
    GraspGen-X  & 71/150 (47.3\%)  & 72 & 62 & 30 \\
    \midrule
    \multicolumn{2}{l}{Hardware success, 150 trials each} \\
    Naive / Iterative / Reflex & \multicolumn{4}{c}{48.0\% / 82.0\% / 86.7\%} \\
    \bottomrule
  \end{tabular}
\end{table}

GraspGen-X is admissible in 47.3\% of trials, level with naive closure and below
iterative and reflex closure. The ranking accounts for much of this: its
top-ranked proposal executes cleanly in none of the 71 trials that succeed at
all, the first collision-free candidate is the 20th proposal by median
(interquartile range 3rd to 24th, at worst the 74th), and 30 of the 71 are found
only beyond the 20th. Scoring the same run over the first 20 proposals, as an
earlier version of this analysis did, gives 27.3\%. Object coverage accounts for
the rest: the generator admits no approach condition at all for the seven objects
whose grasp point is at or below \SI{15}{\milli\metre}, the set on which naive
closure whiffs, and at least 3 of 10 for every object at or above
\SI{29}{\milli\metre} (19 of 50 delicate trials against 52 of 100 YCB-like). It
is also the strategy that admits most trials at the adversarial level-with-object
condition, 5 of 15 against 1 for reflex closure, because a 6-DoF proposal can
approach from the side where a width-conditioned pose cannot.

This is a statement about the interface rather than about the generator's grasp
quality. GraspGen-X is given an object mesh and no support surface, so it
proposes poses from every direction including from beneath the table, and nothing
in its confidence ordering prefers a pose the arm can reach on a tabletop. A
collision check over a deep candidate list recovers part of the gap, and the
depth required is the practical cost: the first 20 proposals recover 27.3\% and
the first 100 recover 47.3\%, and trusting the ranking without such a check
would leave every one of these grasps unexecutable.

\subsection{Limits}
\label{sec:nogo_limits_v2}

The analysis is geometric and bounds what it can explain. On the delicate subset
it places reflex closure above iterative closure, 90.0\% (45/50) against 66.0\%
(33/50), and the hardware benchmark places iterative above reflex, 94\% against
80\%. The
delicate-object failures attributed to reflex closure in
Sec.~\ref{sec:physical_grasp_results} are contact-dynamics events---objects
skidding or being damaged during the faster single-move closure---and no such
mechanism exists in a kinematic model. Every reflex failure here is a support
collision or premature contact. The result is therefore an operating-regime
argument: geometry accounts for the naive failures and for the reflex clearance
advantage at low approaches, and the hold-phase force data carries the
delicate-object case.

Three further limits are quantified rather than asserted. The model omits arm
inverse kinematics. It uses box, cylinder and sphere object proxies rather than
measured meshes. And it assumes a straight-line approach to the pre-grasp pose;
replacing that straight line with an L-shaped travel-then-descend move reproduces
48.0\%, 71.3\% and 90.7\% in all 450 trials, so every support-plane collision
reported here happens during closure rather than during transit. Fingertip
compliance is not modelled at all: adding the silicone pad would move the
fingertip surface outward and make the model more pessimistic, so it cannot
account for the iterative residual.

% ---------------------------------------------------------------------
% APPENDIX MATERIAL
% ---------------------------------------------------------------------

\begin{table}[t]
  \centering
  \caption{Sensitivity of the admissible fraction to the support-plane geometry.
  The capsule tolerances span the measured 5th to 95th percentile of the capsule
  proxy's over-approximation; \SI{10}{\milli\metre} is the value an earlier
  version of this analysis fitted to the hardware rates. The strategy ordering is
  unchanged throughout, and naive closure is insensitive to every setting because
  its failures are terminal misalignments.}
  \label{tab:admissibility_sensitivity}
  \begin{tabular}{lccc}
    \toprule
    Support-plane geometry & Naive & Iterative & Reflex \\
    \midrule
    Capsule proxy, tol.\ \SI{1.8}{\milli\metre}  & 48.0\% & 68.0\% & 90.7\% \\
    Capsule proxy, tol.\ \SI{3.4}{\milli\metre}  & 48.0\% & 68.0\% & 90.7\% \\
    Collision geometry, no tolerance (primary)   & 48.0\% & 71.3\% & 90.7\% \\
    Capsule proxy, tol.\ \SI{6.8}{\milli\metre}  & 48.0\% & 79.3\% & 90.7\% \\
    Capsule proxy, tol.\ \SI{10}{\milli\metre}   & 48.0\% & 84.7\% & 90.7\% \\
    \midrule
    Hardware benchmark (measured success) & 48.0\% & 82.0\% & 86.7\% \\
    \bottomrule
  \end{tabular}
\end{table}

Two model inputs were corrected before the rates above were computed, and their
contributions separate cleanly. Measured object cross-sections replace first-pass
estimates for six objects---and, for the small screwdriver, the handle height is
separated from the \SI{15}{\milli\metre} width it is pinched across, which places
its grasp point at \SI{12}{\milli\metre} rather than \SI{7.5}{\milli\metre} and
is worth 5 iterative and 9 reflex trials---and the approach conditions use the
$\pm\SI{3}{\centi\metre}$ lateral jitter and elevated heights stated in the
hardware protocol in place of a $\pm\SI{150}{\milli\metre}$ reconstruction. With
the estimated sizes and the wide reconstruction the model gives 48.7\%, 65.3\%
and 77.3\%; correcting only the conditions gives 60.0\%, 84.7\% and 90.7\%;
correcting only the sizes gives 42.0\%, 65.3\% and 82.7\%. The conditions carry
iterative closure, because the $\pm\SI{150}{\milli\metre}$ starts require a
diagonal transit that drives the fingers through the support plane on the way in;
the sizes carry naive closure, by moving two more objects below the height at
which its terminal contact misses.

% where Sec. III.C (the current no-go subsection) sits.
%
% It is written to REPLACE the retained proxy-model subsection
% (\label{sec:nogo_analysis}) and its Fig. 7.  Labels carry the suffix
% _v2 so the two can coexist while the swap is made.
%
% Figure 1 below is a drop-in REPLACEMENT for fig:grasp_strategies: it
% keeps img/grasp_strategies_complex.png and adds two no-go strips whose
% three columns line up with the three photographed strategies.
%
% Numbers here are produced by:
%   code/calibrate_capsule_envelope.py --out results/calibration
%   code/run_reachability_sweep.py --sampler paper-conditions \
%       --contact-policy contact-fingers --floor-model mesh \
%       --floor-tolerance-mm 0 --tag paper_conditions_meshfloor
%   code/run_nogo_volume.py --object paper_strawberry --region quarter-ball \
%       --full-grid --floor-model mesh --floor-tolerance-mm 0
%   code/render_strips.sh   (render_nogo_volume.py --strip --panels-dir,
%                            then compose_nogo_figure.py for the lettered grid)
% Calibration and ablations: results/calibration/README.md
% =====================================================================

\subsection{Strategy-Conditioned Approach Feasibility}
\label{sec:nogo_analysis_v2}

A terminal grasp configuration does not specify a collision-free motion to that
configuration. Because the RH56 finger linkages are coupled, a commanded width
also sets fingertip position and the hand pose required to hold a grasp plane, so
the admissible set depends on which closure strategy issues the command. We
model each strategy as an ordered sequence of waypoints in hand-base position,
hand-base rotation, and the six actuator commands, and check the whole commanded
motion---approach and closure---against the object and the support surface.

Let $\mathbf{x}$ be the world hand-base position at which the motion starts and
let $\sigma$ index the strategies. Writing $A(\mathbf{x},\sigma)=1$ when the
commanded motion is admissible, we report
\begin{equation}
  \nu(\sigma) = \frac{1}{|\mathcal{X}|}\sum_{\mathbf{x}\in\mathcal{X}} A(\mathbf{x},\sigma),
  \label{eq:admissible_fraction}
\end{equation}
the fraction of sampled start poses admitted for that strategy. The sample
$\mathcal{X}$ is the ten approach conditions per object used on hardware---one
adversarial pose level with the grasp point, one overhead, and eight at
\SI{150}{\milli\metre} and \SI{250}{\milli\metre} above it with
$\pm\SI{30}{\milli\metre}$ of lateral jitter---over the same 15 objects, giving
150 trials per strategy.

Admissibility requires three geometric conditions. No part of the hand may
penetrate the support plane; no part of the hand outside the planned contact
fingers may penetrate the object; and the antipodal contact midpoint reached at
the end of the motion must lie inside the object. The planned width is the
object width plus \SI{20}{\milli\metre}, so a correctly executed grasp leaves
about \SI{10}{\milli\metre} of clearance per side and any contact in this model
is unintended. The three strategies differ only in what the wrist does during
closure: naive closure holds the wrist at the pose that straddled the object with
the hand open, iterative closure moves the wrist along the analytical width
sequence, and reflex closure holds the wrist at the final analytical pose while
only the opposing fingers move.

The support-plane test is run against the hand's own collision geometry---palm
boxes and cylinder, phalanx meshes---so it carries no clearance tolerance. A
faster 13-capsule proxy of the hand is used for the object tests, and its
vertical over-approximation was measured rather than assumed: over the 41{,}250
hand states the three strategies command, the lowest point of the capsule proxy
lies a median of \SI{3.4}{\milli\metre} below the lowest point of the collision
geometry (5th to 95th percentile \SI{1.8}{}--\SI{6.8}{\milli\metre}), and it lies
below it at every state, so the capsule test is one-sided. The capsule that
reaches lowest is a distal phalanx in 96.6\% of those states. Because MuJoCo
collides a mesh as its convex hull, the same measurement was repeated against the
undecimated visual meshes; the two agree to a median of
\SI{-0.11}{\milli\metre}, so the collision geometry is the part.

\begin{figure}[!htb]
  \centering
  \includegraphics[width=0.95\columnwidth]{img/grasp_strategies_complex.png}\\[3pt]
  % Both objects and all four strategies in one lettered grid; produced by
  % code/render_strips.sh (render_nogo_volume.py --panels-dir, then
  % compose_nogo_figure.py). The three-strategy strips
  % plot/20260908-nogo_strip_{strawberry,can}.png remain available if the
  % GraspGen-X column is dropped.
  \includegraphics[width=0.95\columnwidth]{plot/20260908-nogo_grid_strawberry_can.png}
  \caption{\textbf{Top:} the three grasp strategies---iterative (left), reflex
  (centre), naive (right). In the iterative approach the hand rotates downward
  from a pre-grasp pose to reach the desired grasp pose. In the reflex approach
  the fingers remain open until the thumb makes contact with the object. In the
  naive approach the fingers are treated like a parallel jaw gripper prior to
  closure, leading to ground collisions for small objects.
  \textbf{Bottom:} start poses each strategy can and cannot execute from, one
  column per strategy with GraspGen-X at its best of 10 ranked candidates, for
  (A) the strawberry, grasp point \SI{15}{\milli\metre} above the table, and
  (B) the can, \SI{110}{\milli\metre}. The sampled region is the quarter ball
  of radius \SI{300}{\milli\metre} above the grasp point and on the approach
  side, the poses an arm is commanded from; the red solid is the subset from
  which the strategy cannot execute and the green shell is the rest, so the two
  together are the region each percentage is a fraction of. The object is drawn
  opaque with a post through it. At the strawberry, clearance during closure
  decides the outcome and the strategies separate: reflex closure leaves 53.3\%
  of the region unreachable, iterative closure 76.9\%, and naive closure and
  GraspGen-X all of it. At the can the support surface is never the binding
  constraint and the three hand-designed sets nearly coincide, 34.0\%, 32.1\%
  and 34.0\%, while GraspGen-X leaves 67.7\% unreachable at every sampled
  height, a restriction on approach direction rather than on clearance.
  Searching all 100 proposals leaves the strawberry region unreachable, since
  every proposal places the hand in the table at the proposed pose itself, and
  about 3\% of the can region.
  Surfaces are the half level of the \SI{30}{\milli\metre} sweep grid, refined
  and low-passed, so the smoothness is presentational and the boundary is
  resolved to \SI{30}{\milli\metre}.}
  \label{fig:grasp_strategies_v2}
\end{figure}

Naive closure is admissible in 48.0\% of trials, iterative closure in 71.3\% and
reflex closure in 90.7\%, against measured hardware success rates of 48.0\%,
82.0\% and 86.7\% over 150 trials each. Naive and reflex closure fall inside the
hardware 95\% confidence intervals (39.8--56.3\% and 80.2--91.7\%); iterative
closure is 10.7 points low. No parameter in the model is fitted to a success
rate.

The naive failures separate into two mechanisms, and neither depends on the start
pose. Terminal misalignment accounts for 60 of 150 trials and occurs for no other
hand-designed strategy: holding the wrist fixed through closure leaves the
coupled fingertip travel uncompensated, displacing the contact midpoint by a
median \SI{17.4}{\milli\metre} (\SI{16.3}{}--\SI{18.6}{\milli\metre}), which
carries it outside every object of grasp width \SI{25}{\milli\metre} or less.
Iterative and reflex closure reach their targets to within
\SI{0.004}{\milli\metre} by construction. The result is a step in grasp-point
height rather than a gradient: naive closure admits none of the 70 trials on
objects whose grasp point lies at or below \SI{15}{\milli\metre} and 72 of the 80
at or above \SI{29}{\milli\metre}, and its rate is identical at both elevated
approach heights (53.3\% at \SI{150}{\milli\metre} and at
\SI{250}{\milli\metre}) because the mechanism is set by object geometry.
Support-surface collision is recorded in 76 naive trials, at a median
penetration of \SI{7.9}{\milli\metre}. The two failure modes overlap rather
than add: every one of the 60 misaligned trials also collides with the support
surface, the index intermediate phalanx reaching the plane in 54 of them at a
median depth of \SI{6.1}{\milli\metre}, so the 78 unreachable naive trials are
the 76 support collisions plus 2 premature object contacts with no collision.
The counts in Table~\ref{tab:ggx_admissibility} are of trials showing each mode
and not a partition of the failures.

Reflex closure retains the low-clearance region that iterative closure loses, and
the advantage is confined to the low approach. At the \SI{250}{\milli\metre}
approach height the two are identical at 100.0\%; at \SI{150}{\milli\metre}
reflex closure holds 100.0\% while iterative closure falls to 53.3\%, the same
rate as naive. Reflex closure is also the only hand-designed strategy admitting
any trial at the adversarial condition level with the grasp point (1 of 15,
against 0 and 0). Its 12 support-plane collisions are uniformly deep---the
shallowest is \SI{51.5}{\milli\metre} and the median
\SI{100.9}{\milli\metre}---and correspond to poses no clearance allowance
recovers, whereas 13 of iterative closure's 41 collisions are grazes no deeper
than \SI{1.7}{\milli\metre}.

That graze population is where the 10.7-point iterative residual sits. Forgiving
support-plane contact shallower than \SI{2}{\milli\metre} moves iterative closure
to 79.3\% and shallower than \SI{5}{\milli\metre} to 84.7\%, bracketing the
measured 82.0\%, on the argument that a silicone fingertip deflects at that depth.
Naive and reflex closure do not move, because neither has a shallow collision to
forgive. Testing the argument needs the per-trial failure cause from the hardware
benchmark, which the archive does not record, so the primary configuration
forgives nothing.

\subsection{A Learned Generator Under the Same Test}
\label{sec:ggx_reachability_v2}

The same model accepts a learned grasp generator without modification.
GraspGen-X proposes a 6-DoF hand pose rather than an aperture, and its execution
has the shape of the wrist-fixed strategies: approach the proposed pose with the
hand open, then interpolate the descriptor's closure with the wrist held still.
We run the released model with the official Inspire-hand descriptor over the same
15 object proxies, producing 100 ranked candidates each, and test candidates in
descending confidence order; a start pose is admissible when any of the 100
executes cleanly, and a pose with no clean candidate is a planning failure, which
is the hardware protocol of Sec.~\ref{sec:vision_evaluation}. Its
terminal check cannot be the antipodal-midpoint test, which presumes a precision
pinch, so we use the descriptor's own declared capture volume and ask whether the
proposal engages the object on its terms.

\begin{table}[t]
  \centering
  \caption{All four strategies over 15 objects and ten approach conditions,
  support plane checked against the hand's collision geometry with no clearance
  tolerance. A trial can trigger more than one failure mode.}
  \label{tab:ggx_admissibility}
  \setlength{\tabcolsep}{4pt}
  \begin{tabular}{lcccc}
    \toprule
    & & Support & Premature & Terminal \\
    Strategy & Admissible & collision & contact & miss \\
    \midrule
    Naive       & 72/150 (48.0\%)  & 76 & 10 & 60 \\
    Iterative   & 107/150 (71.3\%) & 41 & 10 &  0 \\
    Reflex      & 136/150 (90.7\%) & 12 &  5 &  0 \\
    GraspGen-X  & 71/150 (47.3\%)  & 72 & 62 & 30 \\
    \midrule
    \multicolumn{2}{l}{Hardware success, 150 trials each} \\
    Naive / Iterative / Reflex & \multicolumn{4}{c}{48.0\% / 82.0\% / 86.7\%} \\
    \bottomrule
  \end{tabular}
\end{table}

GraspGen-X is admissible in 47.3\% of trials, level with naive closure and below
iterative and reflex closure. The ranking accounts for much of this: its
top-ranked proposal executes cleanly in none of the 71 trials that succeed at
all, the first collision-free candidate is the 20th proposal by median
(interquartile range 3rd to 24th, at worst the 74th), and 30 of the 71 are found
only beyond the 20th. Scoring the same run over the first 20 proposals, as an
earlier version of this analysis did, gives 27.3\%. Object coverage accounts for
the rest: the generator admits no approach condition at all for the seven objects
whose grasp point is at or below \SI{15}{\milli\metre}, the set on which naive
closure whiffs, and at least 3 of 10 for every object at or above
\SI{29}{\milli\metre} (19 of 50 delicate trials against 52 of 100 YCB-like). It
is also the strategy that admits most trials at the adversarial level-with-object
condition, 5 of 15 against 1 for reflex closure, because a 6-DoF proposal can
approach from the side where a width-conditioned pose cannot.

This is a statement about the interface rather than about the generator's grasp
quality. GraspGen-X is given an object mesh and no support surface, so it
proposes poses from every direction including from beneath the table, and nothing
in its confidence ordering prefers a pose the arm can reach on a tabletop. A
collision check over a deep candidate list recovers part of the gap, and the
depth required is the practical cost: the first 20 proposals recover 27.3\% and
the first 100 recover 47.3\%, and trusting the ranking without such a check
would leave every one of these grasps unexecutable.

\subsection{Limits}
\label{sec:nogo_limits_v2}

The analysis is geometric and bounds what it can explain. On the delicate subset
it places reflex closure above iterative closure, 90.0\% (45/50) against 66.0\%
(33/50), and the hardware benchmark places iterative above reflex, 94\% against
80\%. The
delicate-object failures attributed to reflex closure in
Sec.~\ref{sec:physical_grasp_results} are contact-dynamics events---objects
skidding or being damaged during the faster single-move closure---and no such
mechanism exists in a kinematic model. Every reflex failure here is a support
collision or premature contact. The result is therefore an operating-regime
argument: geometry accounts for the naive failures and for the reflex clearance
advantage at low approaches, and the hold-phase force data carries the
delicate-object case.

Three further limits are quantified rather than asserted. The model omits arm
inverse kinematics. It uses box, cylinder and sphere object proxies rather than
measured meshes. And it assumes a straight-line approach to the pre-grasp pose;
replacing that straight line with an L-shaped travel-then-descend move reproduces
48.0\%, 71.3\% and 90.7\% in all 450 trials, so every support-plane collision
reported here happens during closure rather than during transit. Fingertip
compliance is not modelled at all: adding the silicone pad would move the
fingertip surface outward and make the model more pessimistic, so it cannot
account for the iterative residual.

% ---------------------------------------------------------------------
% APPENDIX MATERIAL
% ---------------------------------------------------------------------

\begin{table}[t]
  \centering
  \caption{Sensitivity of the admissible fraction to the support-plane geometry.
  The capsule tolerances span the measured 5th to 95th percentile of the capsule
  proxy's over-approximation; \SI{10}{\milli\metre} is the value an earlier
  version of this analysis fitted to the hardware rates. The strategy ordering is
  unchanged throughout, and naive closure is insensitive to every setting because
  its failures are terminal misalignments.}
  \label{tab:admissibility_sensitivity}
  \begin{tabular}{lccc}
    \toprule
    Support-plane geometry & Naive & Iterative & Reflex \\
    \midrule
    Capsule proxy, tol.\ \SI{1.8}{\milli\metre}  & 48.0\% & 68.0\% & 90.7\% \\
    Capsule proxy, tol.\ \SI{3.4}{\milli\metre}  & 48.0\% & 68.0\% & 90.7\% \\
    Collision geometry, no tolerance (primary)   & 48.0\% & 71.3\% & 90.7\% \\
    Capsule proxy, tol.\ \SI{6.8}{\milli\metre}  & 48.0\% & 79.3\% & 90.7\% \\
    Capsule proxy, tol.\ \SI{10}{\milli\metre}   & 48.0\% & 84.7\% & 90.7\% \\
    \midrule
    Hardware benchmark (measured success) & 48.0\% & 82.0\% & 86.7\% \\
    \bottomrule
  \end{tabular}
\end{table}

Two model inputs were corrected before the rates above were computed, and their
contributions separate cleanly. Measured object cross-sections replace first-pass
estimates for six objects---and, for the small screwdriver, the handle height is
separated from the \SI{15}{\milli\metre} width it is pinched across, which places
its grasp point at \SI{12}{\milli\metre} rather than \SI{7.5}{\milli\metre} and
is worth 5 iterative and 9 reflex trials---and the approach conditions use the
$\pm\SI{3}{\centi\metre}$ lateral jitter and elevated heights stated in the
hardware protocol in place of a $\pm\SI{150}{\milli\metre}$ reconstruction. With
the estimated sizes and the wide reconstruction the model gives 48.7\%, 65.3\%
and 77.3\%; correcting only the conditions gives 60.0\%, 84.7\% and 90.7\%;
correcting only the sizes gives 42.0\%, 65.3\% and 82.7\%. The conditions carry
iterative closure, because the $\pm\SI{150}{\milli\metre}$ starts require a
diagonal transit that drives the fingers through the support plane on the way in;
the sizes carry naive closure, by moving two more objects below the height at
which its terminal contact misses.

% where Sec. III.C (the current no-go subsection) sits.
%
% It is written to REPLACE the retained proxy-model subsection
% (\label{sec:nogo_analysis}) and its Fig. 7.  Labels carry the suffix
% _v2 so the two can coexist while the swap is made.
%
% Figure 1 below is a drop-in REPLACEMENT for fig:grasp_strategies: it
% keeps img/grasp_strategies_complex.png and adds two no-go strips whose
% three columns line up with the three photographed strategies.
%
% Numbers here are produced by:
%   code/calibrate_capsule_envelope.py --out results/calibration
%   code/run_reachability_sweep.py --sampler paper-conditions \
%       --contact-policy contact-fingers --floor-model mesh \
%       --floor-tolerance-mm 0 --tag paper_conditions_meshfloor
%   code/run_nogo_volume.py --object paper_strawberry --region quarter-ball \
%       --full-grid --floor-model mesh --floor-tolerance-mm 0
%   code/render_strips.sh   (render_nogo_volume.py --strip --panels-dir,
%                            then compose_nogo_figure.py for the lettered grid)
% Calibration and ablations: results/calibration/README.md
% =====================================================================

\subsection{Strategy-Conditioned Approach Feasibility}
\label{sec:nogo_analysis_v2}

A terminal grasp configuration does not specify a collision-free motion to that
configuration. Because the RH56 finger linkages are coupled, a commanded width
also sets fingertip position and the hand pose required to hold a grasp plane, so
the admissible set depends on which closure strategy issues the command. We
model each strategy as an ordered sequence of waypoints in hand-base position,
hand-base rotation, and the six actuator commands, and check the whole commanded
motion---approach and closure---against the object and the support surface.

Let $\mathbf{x}$ be the world hand-base position at which the motion starts and
let $\sigma$ index the strategies. Writing $A(\mathbf{x},\sigma)=1$ when the
commanded motion is admissible, we report
\begin{equation}
  \nu(\sigma) = \frac{1}{|\mathcal{X}|}\sum_{\mathbf{x}\in\mathcal{X}} A(\mathbf{x},\sigma),
  \label{eq:admissible_fraction}
\end{equation}
the fraction of sampled start poses admitted for that strategy. The sample
$\mathcal{X}$ is the ten approach conditions per object used on hardware---one
adversarial pose level with the grasp point, one overhead, and eight at
\SI{150}{\milli\metre} and \SI{250}{\milli\metre} above it with
$\pm\SI{30}{\milli\metre}$ of lateral jitter---over the same 15 objects, giving
150 trials per strategy.

Admissibility requires three geometric conditions. No part of the hand may
penetrate the support plane; no part of the hand outside the planned contact
fingers may penetrate the object; and the antipodal contact midpoint reached at
the end of the motion must lie inside the object. The planned width is the
object width plus \SI{20}{\milli\metre}, so a correctly executed grasp leaves
about \SI{10}{\milli\metre} of clearance per side and any contact in this model
is unintended. The three strategies differ only in what the wrist does during
closure: naive closure holds the wrist at the pose that straddled the object with
the hand open, iterative closure moves the wrist along the analytical width
sequence, and reflex closure holds the wrist at the final analytical pose while
only the opposing fingers move.

The support-plane test is run against the hand's own collision geometry---palm
boxes and cylinder, phalanx meshes---so it carries no clearance tolerance. A
faster 13-capsule proxy of the hand is used for the object tests, and its
vertical over-approximation was measured rather than assumed: over the 41{,}250
hand states the three strategies command, the lowest point of the capsule proxy
lies a median of \SI{3.4}{\milli\metre} below the lowest point of the collision
geometry (5th to 95th percentile \SI{1.8}{}--\SI{6.8}{\milli\metre}), and it lies
below it at every state, so the capsule test is one-sided. The capsule that
reaches lowest is a distal phalanx in 96.6\% of those states. Because MuJoCo
collides a mesh as its convex hull, the same measurement was repeated against the
undecimated visual meshes; the two agree to a median of
\SI{-0.11}{\milli\metre}, so the collision geometry is the part.

\begin{figure}[!htb]
  \centering
  \includegraphics[width=0.95\columnwidth]{img/grasp_strategies_complex.png}\\[3pt]
  % Both objects and all four strategies in one lettered grid; produced by
  % code/render_strips.sh (render_nogo_volume.py --panels-dir, then
  % compose_nogo_figure.py). The three-strategy strips
  % plot/20260908-nogo_strip_{strawberry,can}.png remain available if the
  % GraspGen-X column is dropped.
  \includegraphics[width=0.95\columnwidth]{plot/20260908-nogo_grid_strawberry_can.png}
  \caption{\textbf{Top:} the three grasp strategies---iterative (left), reflex
  (centre), naive (right). In the iterative approach the hand rotates downward
  from a pre-grasp pose to reach the desired grasp pose. In the reflex approach
  the fingers remain open until the thumb makes contact with the object. In the
  naive approach the fingers are treated like a parallel jaw gripper prior to
  closure, leading to ground collisions for small objects.
  \textbf{Bottom:} start poses each strategy can and cannot execute from, one
  column per strategy with GraspGen-X at its best of 10 ranked candidates, for
  (A) the strawberry, grasp point \SI{15}{\milli\metre} above the table, and
  (B) the can, \SI{110}{\milli\metre}. The sampled region is the quarter ball
  of radius \SI{300}{\milli\metre} above the grasp point and on the approach
  side, the poses an arm is commanded from; the red solid is the subset from
  which the strategy cannot execute and the green shell is the rest, so the two
  together are the region each percentage is a fraction of. The object is drawn
  opaque with a post through it. At the strawberry, clearance during closure
  decides the outcome and the strategies separate: reflex closure leaves 53.3\%
  of the region unreachable, iterative closure 76.9\%, and naive closure and
  GraspGen-X all of it. At the can the support surface is never the binding
  constraint and the three hand-designed sets nearly coincide, 34.0\%, 32.1\%
  and 34.0\%, while GraspGen-X leaves 67.7\% unreachable at every sampled
  height, a restriction on approach direction rather than on clearance.
  Searching all 100 proposals leaves the strawberry region unreachable, since
  every proposal places the hand in the table at the proposed pose itself, and
  about 3\% of the can region.
  Surfaces are the half level of the \SI{30}{\milli\metre} sweep grid, refined
  and low-passed, so the smoothness is presentational and the boundary is
  resolved to \SI{30}{\milli\metre}.}
  \label{fig:grasp_strategies_v2}
\end{figure}

Naive closure is admissible in 48.0\% of trials, iterative closure in 71.3\% and
reflex closure in 90.7\%, against measured hardware success rates of 48.0\%,
82.0\% and 86.7\% over 150 trials each. Naive and reflex closure fall inside the
hardware 95\% confidence intervals (39.8--56.3\% and 80.2--91.7\%); iterative
closure is 10.7 points low. No parameter in the model is fitted to a success
rate.

The naive failures separate into two mechanisms, and neither depends on the start
pose. Terminal misalignment accounts for 60 of 150 trials and occurs for no other
hand-designed strategy: holding the wrist fixed through closure leaves the
coupled fingertip travel uncompensated, displacing the contact midpoint by a
median \SI{17.4}{\milli\metre} (\SI{16.3}{}--\SI{18.6}{\milli\metre}), which
carries it outside every object of grasp width \SI{25}{\milli\metre} or less.
Iterative and reflex closure reach their targets to within
\SI{0.004}{\milli\metre} by construction. The result is a step in grasp-point
height rather than a gradient: naive closure admits none of the 70 trials on
objects whose grasp point lies at or below \SI{15}{\milli\metre} and 72 of the 80
at or above \SI{29}{\milli\metre}, and its rate is identical at both elevated
approach heights (53.3\% at \SI{150}{\milli\metre} and at
\SI{250}{\milli\metre}) because the mechanism is set by object geometry.
Support-surface collision is recorded in 76 naive trials, at a median
penetration of \SI{7.9}{\milli\metre}. The two failure modes overlap rather
than add: every one of the 60 misaligned trials also collides with the support
surface, the index intermediate phalanx reaching the plane in 54 of them at a
median depth of \SI{6.1}{\milli\metre}, so the 78 unreachable naive trials are
the 76 support collisions plus 2 premature object contacts with no collision.
The counts in Table~\ref{tab:ggx_admissibility} are of trials showing each mode
and not a partition of the failures.

Reflex closure retains the low-clearance region that iterative closure loses, and
the advantage is confined to the low approach. At the \SI{250}{\milli\metre}
approach height the two are identical at 100.0\%; at \SI{150}{\milli\metre}
reflex closure holds 100.0\% while iterative closure falls to 53.3\%, the same
rate as naive. Reflex closure is also the only hand-designed strategy admitting
any trial at the adversarial condition level with the grasp point (1 of 15,
against 0 and 0). Its 12 support-plane collisions are uniformly deep---the
shallowest is \SI{51.5}{\milli\metre} and the median
\SI{100.9}{\milli\metre}---and correspond to poses no clearance allowance
recovers, whereas 13 of iterative closure's 41 collisions are grazes no deeper
than \SI{1.7}{\milli\metre}.

That graze population is where the 10.7-point iterative residual sits. Forgiving
support-plane contact shallower than \SI{2}{\milli\metre} moves iterative closure
to 79.3\% and shallower than \SI{5}{\milli\metre} to 84.7\%, bracketing the
measured 82.0\%, on the argument that a silicone fingertip deflects at that depth.
Naive and reflex closure do not move, because neither has a shallow collision to
forgive. Testing the argument needs the per-trial failure cause from the hardware
benchmark, which the archive does not record, so the primary configuration
forgives nothing.

\subsection{A Learned Generator Under the Same Test}
\label{sec:ggx_reachability_v2}

The same model accepts a learned grasp generator without modification.
GraspGen-X proposes a 6-DoF hand pose rather than an aperture, and its execution
has the shape of the wrist-fixed strategies: approach the proposed pose with the
hand open, then interpolate the descriptor's closure with the wrist held still.
We run the released model with the official Inspire-hand descriptor over the same
15 object proxies, producing 100 ranked candidates each, and test candidates in
descending confidence order; a start pose is admissible when any of the 100
executes cleanly, and a pose with no clean candidate is a planning failure, which
is the hardware protocol of Sec.~\ref{sec:vision_evaluation}. Its
terminal check cannot be the antipodal-midpoint test, which presumes a precision
pinch, so we use the descriptor's own declared capture volume and ask whether the
proposal engages the object on its terms.

\begin{table}[t]
  \centering
  \caption{All four strategies over 15 objects and ten approach conditions,
  support plane checked against the hand's collision geometry with no clearance
  tolerance. A trial can trigger more than one failure mode.}
  \label{tab:ggx_admissibility}
  \setlength{\tabcolsep}{4pt}
  \begin{tabular}{lcccc}
    \toprule
    & & Support & Premature & Terminal \\
    Strategy & Admissible & collision & contact & miss \\
    \midrule
    Naive       & 72/150 (48.0\%)  & 76 & 10 & 60 \\
    Iterative   & 107/150 (71.3\%) & 41 & 10 &  0 \\
    Reflex      & 136/150 (90.7\%) & 12 &  5 &  0 \\
    GraspGen-X  & 71/150 (47.3\%)  & 72 & 62 & 30 \\
    \midrule
    \multicolumn{2}{l}{Hardware success, 150 trials each} \\
    Naive / Iterative / Reflex & \multicolumn{4}{c}{48.0\% / 82.0\% / 86.7\%} \\
    \bottomrule
  \end{tabular}
\end{table}

GraspGen-X is admissible in 47.3\% of trials, level with naive closure and below
iterative and reflex closure. The ranking accounts for much of this: its
top-ranked proposal executes cleanly in none of the 71 trials that succeed at
all, the first collision-free candidate is the 20th proposal by median
(interquartile range 3rd to 24th, at worst the 74th), and 30 of the 71 are found
only beyond the 20th. Scoring the same run over the first 20 proposals, as an
earlier version of this analysis did, gives 27.3\%. Object coverage accounts for
the rest: the generator admits no approach condition at all for the seven objects
whose grasp point is at or below \SI{15}{\milli\metre}, the set on which naive
closure whiffs, and at least 3 of 10 for every object at or above
\SI{29}{\milli\metre} (19 of 50 delicate trials against 52 of 100 YCB-like). It
is also the strategy that admits most trials at the adversarial level-with-object
condition, 5 of 15 against 1 for reflex closure, because a 6-DoF proposal can
approach from the side where a width-conditioned pose cannot.

This is a statement about the interface rather than about the generator's grasp
quality. GraspGen-X is given an object mesh and no support surface, so it
proposes poses from every direction including from beneath the table, and nothing
in its confidence ordering prefers a pose the arm can reach on a tabletop. A
collision check over a deep candidate list recovers part of the gap, and the
depth required is the practical cost: the first 20 proposals recover 27.3\% and
the first 100 recover 47.3\%, and trusting the ranking without such a check
would leave every one of these grasps unexecutable.

\subsection{Limits}
\label{sec:nogo_limits_v2}

The analysis is geometric and bounds what it can explain. On the delicate subset
it places reflex closure above iterative closure, 90.0\% (45/50) against 66.0\%
(33/50), and the hardware benchmark places iterative above reflex, 94\% against
80\%. The
delicate-object failures attributed to reflex closure in
Sec.~\ref{sec:physical_grasp_results} are contact-dynamics events---objects
skidding or being damaged during the faster single-move closure---and no such
mechanism exists in a kinematic model. Every reflex failure here is a support
collision or premature contact. The result is therefore an operating-regime
argument: geometry accounts for the naive failures and for the reflex clearance
advantage at low approaches, and the hold-phase force data carries the
delicate-object case.

Three further limits are quantified rather than asserted. The model omits arm
inverse kinematics. It uses box, cylinder and sphere object proxies rather than
measured meshes. And it assumes a straight-line approach to the pre-grasp pose;
replacing that straight line with an L-shaped travel-then-descend move reproduces
48.0\%, 71.3\% and 90.7\% in all 450 trials, so every support-plane collision
reported here happens during closure rather than during transit. Fingertip
compliance is not modelled at all: adding the silicone pad would move the
fingertip surface outward and make the model more pessimistic, so it cannot
account for the iterative residual.

% ---------------------------------------------------------------------
% APPENDIX MATERIAL
% ---------------------------------------------------------------------

\begin{table}[t]
  \centering
  \caption{Sensitivity of the admissible fraction to the support-plane geometry.
  The capsule tolerances span the measured 5th to 95th percentile of the capsule
  proxy's over-approximation; \SI{10}{\milli\metre} is the value an earlier
  version of this analysis fitted to the hardware rates. The strategy ordering is
  unchanged throughout, and naive closure is insensitive to every setting because
  its failures are terminal misalignments.}
  \label{tab:admissibility_sensitivity}
  \begin{tabular}{lccc}
    \toprule
    Support-plane geometry & Naive & Iterative & Reflex \\
    \midrule
    Capsule proxy, tol.\ \SI{1.8}{\milli\metre}  & 48.0\% & 68.0\% & 90.7\% \\
    Capsule proxy, tol.\ \SI{3.4}{\milli\metre}  & 48.0\% & 68.0\% & 90.7\% \\
    Collision geometry, no tolerance (primary)   & 48.0\% & 71.3\% & 90.7\% \\
    Capsule proxy, tol.\ \SI{6.8}{\milli\metre}  & 48.0\% & 79.3\% & 90.7\% \\
    Capsule proxy, tol.\ \SI{10}{\milli\metre}   & 48.0\% & 84.7\% & 90.7\% \\
    \midrule
    Hardware benchmark (measured success) & 48.0\% & 82.0\% & 86.7\% \\
    \bottomrule
  \end{tabular}
\end{table}

Two model inputs were corrected before the rates above were computed, and their
contributions separate cleanly. Measured object cross-sections replace first-pass
estimates for six objects---and, for the small screwdriver, the handle height is
separated from the \SI{15}{\milli\metre} width it is pinched across, which places
its grasp point at \SI{12}{\milli\metre} rather than \SI{7.5}{\milli\metre} and
is worth 5 iterative and 9 reflex trials---and the approach conditions use the
$\pm\SI{3}{\centi\metre}$ lateral jitter and elevated heights stated in the
hardware protocol in place of a $\pm\SI{150}{\milli\metre}$ reconstruction. With
the estimated sizes and the wide reconstruction the model gives 48.7\%, 65.3\%
and 77.3\%; correcting only the conditions gives 60.0\%, 84.7\% and 90.7\%;
correcting only the sizes gives 42.0\%, 65.3\% and 82.7\%. The conditions carry
iterative closure, because the $\pm\SI{150}{\milli\metre}$ starts require a
diagonal transit that drives the fingers through the support plane on the way in;
the sizes carry naive closure, by moving two more objects below the height at
which its terminal contact misses.

% where Sec. III.C (the current no-go subsection) sits.
%
% It is written to REPLACE the retained proxy-model subsection
% (\label{sec:nogo_analysis}) and its Fig. 7.  Labels carry the suffix
% _v2 so the two can coexist while the swap is made.
%
% Figure 1 below is a drop-in REPLACEMENT for fig:grasp_strategies: it
% keeps img/grasp_strategies_complex.png and adds two no-go strips whose
% three columns line up with the three photographed strategies.
%
% Numbers here are produced by:
%   code/calibrate_capsule_envelope.py --out results/calibration
%   code/run_reachability_sweep.py --sampler paper-conditions \
%       --contact-policy contact-fingers --floor-model mesh \
%       --floor-tolerance-mm 0 --tag paper_conditions_meshfloor
%   code/run_nogo_volume.py --object paper_strawberry --region quarter-ball \
%       --full-grid --floor-model mesh --floor-tolerance-mm 0
%   code/render_strips.sh   (render_nogo_volume.py --strip --panels-dir,
%                            then compose_nogo_figure.py for the lettered grid)
% Calibration and ablations: results/calibration/README.md
% =====================================================================

\subsection{Strategy-Conditioned Approach Feasibility}
\label{sec:nogo_analysis_v2}

A terminal grasp configuration does not specify a collision-free motion to that
configuration. Because the RH56 finger linkages are coupled, a commanded width
also sets fingertip position and the hand pose required to hold a grasp plane, so
the admissible set depends on which closure strategy issues the command. We
model each strategy as an ordered sequence of waypoints in hand-base position,
hand-base rotation, and the six actuator commands, and check the whole commanded
motion---approach and closure---against the object and the support surface.

Let $\mathbf{x}$ be the world hand-base position at which the motion starts and
let $\sigma$ index the strategies. Writing $A(\mathbf{x},\sigma)=1$ when the
commanded motion is admissible, we report
\begin{equation}
  \nu(\sigma) = \frac{1}{|\mathcal{X}|}\sum_{\mathbf{x}\in\mathcal{X}} A(\mathbf{x},\sigma),
  \label{eq:admissible_fraction}
\end{equation}
the fraction of sampled start poses admitted for that strategy. The sample
$\mathcal{X}$ is the ten approach conditions per object used on hardware---one
adversarial pose level with the grasp point, one overhead, and eight at
\SI{150}{\milli\metre} and \SI{250}{\milli\metre} above it with
$\pm\SI{30}{\milli\metre}$ of lateral jitter---over the same 15 objects, giving
150 trials per strategy.

Admissibility requires three geometric conditions. No part of the hand may
penetrate the support plane; no part of the hand outside the planned contact
fingers may penetrate the object; and the antipodal contact midpoint reached at
the end of the motion must lie inside the object. The planned width is the
object width plus \SI{20}{\milli\metre}, so a correctly executed grasp leaves
about \SI{10}{\milli\metre} of clearance per side and any contact in this model
is unintended. The three strategies differ only in what the wrist does during
closure: naive closure holds the wrist at the pose that straddled the object with
the hand open, iterative closure moves the wrist along the analytical width
sequence, and reflex closure holds the wrist at the final analytical pose while
only the opposing fingers move.

The support-plane test is run against the hand's own collision geometry---palm
boxes and cylinder, phalanx meshes---so it carries no clearance tolerance. A
faster 13-capsule proxy of the hand is used for the object tests, and its
vertical over-approximation was measured rather than assumed: over the 41{,}250
hand states the three strategies command, the lowest point of the capsule proxy
lies a median of \SI{3.4}{\milli\metre} below the lowest point of the collision
geometry (5th to 95th percentile \SI{1.8}{}--\SI{6.8}{\milli\metre}), and it lies
below it at every state, so the capsule test is one-sided. The capsule that
reaches lowest is a distal phalanx in 96.6\% of those states. Because MuJoCo
collides a mesh as its convex hull, the same measurement was repeated against the
undecimated visual meshes; the two agree to a median of
\SI{-0.11}{\milli\metre}, so the collision geometry is the part.

\begin{figure}[!htb]
  \centering
  \includegraphics[width=0.95\columnwidth]{img/grasp_strategies_complex.png}\\[3pt]
  % Both objects and all four strategies in one lettered grid; produced by
  % code/render_strips.sh (render_nogo_volume.py --panels-dir, then
  % compose_nogo_figure.py). The three-strategy strips
  % plot/20260908-nogo_strip_{strawberry,can}.png remain available if the
  % GraspGen-X column is dropped.
  \includegraphics[width=0.95\columnwidth]{plot/20260908-nogo_grid_strawberry_can.png}
  \caption{\textbf{Top:} the three grasp strategies---iterative (left), reflex
  (centre), naive (right). In the iterative approach the hand rotates downward
  from a pre-grasp pose to reach the desired grasp pose. In the reflex approach
  the fingers remain open until the thumb makes contact with the object. In the
  naive approach the fingers are treated like a parallel jaw gripper prior to
  closure, leading to ground collisions for small objects.
  \textbf{Bottom:} start poses each strategy can and cannot execute from, one
  column per strategy with GraspGen-X at its best of 10 ranked candidates, for
  (A) the strawberry, grasp point \SI{15}{\milli\metre} above the table, and
  (B) the can, \SI{110}{\milli\metre}. The sampled region is the quarter ball
  of radius \SI{300}{\milli\metre} above the grasp point and on the approach
  side, the poses an arm is commanded from; the red solid is the subset from
  which the strategy cannot execute and the green shell is the rest, so the two
  together are the region each percentage is a fraction of. The object is drawn
  opaque with a post through it. At the strawberry, clearance during closure
  decides the outcome and the strategies separate: reflex closure leaves 53.3\%
  of the region unreachable, iterative closure 76.9\%, and naive closure and
  GraspGen-X all of it. At the can the support surface is never the binding
  constraint and the three hand-designed sets nearly coincide, 34.0\%, 32.1\%
  and 34.0\%, while GraspGen-X leaves 67.7\% unreachable at every sampled
  height, a restriction on approach direction rather than on clearance.
  Searching all 100 proposals leaves the strawberry region unreachable, since
  every proposal places the hand in the table at the proposed pose itself, and
  about 3\% of the can region.
  Surfaces are the half level of the \SI{30}{\milli\metre} sweep grid, refined
  and low-passed, so the smoothness is presentational and the boundary is
  resolved to \SI{30}{\milli\metre}.}
  \label{fig:grasp_strategies_v2}
\end{figure}

Naive closure is admissible in 48.0\% of trials, iterative closure in 71.3\% and
reflex closure in 90.7\%, against measured hardware success rates of 48.0\%,
82.0\% and 86.7\% over 150 trials each. Naive and reflex closure fall inside the
hardware 95\% confidence intervals (39.8--56.3\% and 80.2--91.7\%); iterative
closure is 10.7 points low. No parameter in the model is fitted to a success
rate.

The naive failures separate into two mechanisms, and neither depends on the start
pose. Terminal misalignment accounts for 60 of 150 trials and occurs for no other
hand-designed strategy: holding the wrist fixed through closure leaves the
coupled fingertip travel uncompensated, displacing the contact midpoint by a
median \SI{17.4}{\milli\metre} (\SI{16.3}{}--\SI{18.6}{\milli\metre}), which
carries it outside every object of grasp width \SI{25}{\milli\metre} or less.
Iterative and reflex closure reach their targets to within
\SI{0.004}{\milli\metre} by construction. The result is a step in grasp-point
height rather than a gradient: naive closure admits none of the 70 trials on
objects whose grasp point lies at or below \SI{15}{\milli\metre} and 72 of the 80
at or above \SI{29}{\milli\metre}, and its rate is identical at both elevated
approach heights (53.3\% at \SI{150}{\milli\metre} and at
\SI{250}{\milli\metre}) because the mechanism is set by object geometry.
Support-surface collision is recorded in 76 naive trials, at a median
penetration of \SI{7.9}{\milli\metre}. The two failure modes overlap rather
than add: every one of the 60 misaligned trials also collides with the support
surface, the index intermediate phalanx reaching the plane in 54 of them at a
median depth of \SI{6.1}{\milli\metre}, so the 78 unreachable naive trials are
the 76 support collisions plus 2 premature object contacts with no collision.
The counts in Table~\ref{tab:ggx_admissibility} are of trials showing each mode
and not a partition of the failures.

Reflex closure retains the low-clearance region that iterative closure loses, and
the advantage is confined to the low approach. At the \SI{250}{\milli\metre}
approach height the two are identical at 100.0\%; at \SI{150}{\milli\metre}
reflex closure holds 100.0\% while iterative closure falls to 53.3\%, the same
rate as naive. Reflex closure is also the only hand-designed strategy admitting
any trial at the adversarial condition level with the grasp point (1 of 15,
against 0 and 0). Its 12 support-plane collisions are uniformly deep---the
shallowest is \SI{51.5}{\milli\metre} and the median
\SI{100.9}{\milli\metre}---and correspond to poses no clearance allowance
recovers, whereas 13 of iterative closure's 41 collisions are grazes no deeper
than \SI{1.7}{\milli\metre}.

That graze population is where the 10.7-point iterative residual sits. Forgiving
support-plane contact shallower than \SI{2}{\milli\metre} moves iterative closure
to 79.3\% and shallower than \SI{5}{\milli\metre} to 84.7\%, bracketing the
measured 82.0\%, on the argument that a silicone fingertip deflects at that depth.
Naive and reflex closure do not move, because neither has a shallow collision to
forgive. Testing the argument needs the per-trial failure cause from the hardware
benchmark, which the archive does not record, so the primary configuration
forgives nothing.

\subsection{A Learned Generator Under the Same Test}
\label{sec:ggx_reachability_v2}

The same model accepts a learned grasp generator without modification.
GraspGen-X proposes a 6-DoF hand pose rather than an aperture, and its execution
has the shape of the wrist-fixed strategies: approach the proposed pose with the
hand open, then interpolate the descriptor's closure with the wrist held still.
We run the released model with the official Inspire-hand descriptor over the same
15 object proxies, producing 100 ranked candidates each, and test candidates in
descending confidence order; a start pose is admissible when any of the 100
executes cleanly, and a pose with no clean candidate is a planning failure, which
is the hardware protocol of Sec.~\ref{sec:vision_evaluation}. Its
terminal check cannot be the antipodal-midpoint test, which presumes a precision
pinch, so we use the descriptor's own declared capture volume and ask whether the
proposal engages the object on its terms.

\begin{table}[t]
  \centering
  \caption{All four strategies over 15 objects and ten approach conditions,
  support plane checked against the hand's collision geometry with no clearance
  tolerance. A trial can trigger more than one failure mode.}
  \label{tab:ggx_admissibility}
  \setlength{\tabcolsep}{4pt}
  \begin{tabular}{lcccc}
    \toprule
    & & Support & Premature & Terminal \\
    Strategy & Admissible & collision & contact & miss \\
    \midrule
    Naive       & 72/150 (48.0\%)  & 76 & 10 & 60 \\
    Iterative   & 107/150 (71.3\%) & 41 & 10 &  0 \\
    Reflex      & 136/150 (90.7\%) & 12 &  5 &  0 \\
    GraspGen-X  & 71/150 (47.3\%)  & 72 & 62 & 30 \\
    \midrule
    \multicolumn{2}{l}{Hardware success, 150 trials each} \\
    Naive / Iterative / Reflex & \multicolumn{4}{c}{48.0\% / 82.0\% / 86.7\%} \\
    \bottomrule
  \end{tabular}
\end{table}

GraspGen-X is admissible in 47.3\% of trials, level with naive closure and below
iterative and reflex closure. The ranking accounts for much of this: its
top-ranked proposal executes cleanly in none of the 71 trials that succeed at
all, the first collision-free candidate is the 20th proposal by median
(interquartile range 3rd to 24th, at worst the 74th), and 30 of the 71 are found
only beyond the 20th. Scoring the same run over the first 20 proposals, as an
earlier version of this analysis did, gives 27.3\%. Object coverage accounts for
the rest: the generator admits no approach condition at all for the seven objects
whose grasp point is at or below \SI{15}{\milli\metre}, the set on which naive
closure whiffs, and at least 3 of 10 for every object at or above
\SI{29}{\milli\metre} (19 of 50 delicate trials against 52 of 100 YCB-like). It
is also the strategy that admits most trials at the adversarial level-with-object
condition, 5 of 15 against 1 for reflex closure, because a 6-DoF proposal can
approach from the side where a width-conditioned pose cannot.

This is a statement about the interface rather than about the generator's grasp
quality. GraspGen-X is given an object mesh and no support surface, so it
proposes poses from every direction including from beneath the table, and nothing
in its confidence ordering prefers a pose the arm can reach on a tabletop. A
collision check over a deep candidate list recovers part of the gap, and the
depth required is the practical cost: the first 20 proposals recover 27.3\% and
the first 100 recover 47.3\%, and trusting the ranking without such a check
would leave every one of these grasps unexecutable.

\subsection{Limits}
\label{sec:nogo_limits_v2}

The analysis is geometric and bounds what it can explain. On the delicate subset
it places reflex closure above iterative closure, 90.0\% (45/50) against 66.0\%
(33/50), and the hardware benchmark places iterative above reflex, 94\% against
80\%. The
delicate-object failures attributed to reflex closure in
Sec.~\ref{sec:physical_grasp_results} are contact-dynamics events---objects
skidding or being damaged during the faster single-move closure---and no such
mechanism exists in a kinematic model. Every reflex failure here is a support
collision or premature contact. The result is therefore an operating-regime
argument: geometry accounts for the naive failures and for the reflex clearance
advantage at low approaches, and the hold-phase force data carries the
delicate-object case.

Three further limits are quantified rather than asserted. The model omits arm
inverse kinematics. It uses box, cylinder and sphere object proxies rather than
measured meshes. And it assumes a straight-line approach to the pre-grasp pose;
replacing that straight line with an L-shaped travel-then-descend move reproduces
48.0\%, 71.3\% and 90.7\% in all 450 trials, so every support-plane collision
reported here happens during closure rather than during transit. Fingertip
compliance is not modelled at all: adding the silicone pad would move the
fingertip surface outward and make the model more pessimistic, so it cannot
account for the iterative residual.

% ---------------------------------------------------------------------
% APPENDIX MATERIAL
% ---------------------------------------------------------------------

\begin{table}[t]
  \centering
  \caption{Sensitivity of the admissible fraction to the support-plane geometry.
  The capsule tolerances span the measured 5th to 95th percentile of the capsule
  proxy's over-approximation; \SI{10}{\milli\metre} is the value an earlier
  version of this analysis fitted to the hardware rates. The strategy ordering is
  unchanged throughout, and naive closure is insensitive to every setting because
  its failures are terminal misalignments.}
  \label{tab:admissibility_sensitivity}
  \begin{tabular}{lccc}
    \toprule
    Support-plane geometry & Naive & Iterative & Reflex \\
    \midrule
    Capsule proxy, tol.\ \SI{1.8}{\milli\metre}  & 48.0\% & 68.0\% & 90.7\% \\
    Capsule proxy, tol.\ \SI{3.4}{\milli\metre}  & 48.0\% & 68.0\% & 90.7\% \\
    Collision geometry, no tolerance (primary)   & 48.0\% & 71.3\% & 90.7\% \\
    Capsule proxy, tol.\ \SI{6.8}{\milli\metre}  & 48.0\% & 79.3\% & 90.7\% \\
    Capsule proxy, tol.\ \SI{10}{\milli\metre}   & 48.0\% & 84.7\% & 90.7\% \\
    \midrule
    Hardware benchmark (measured success) & 48.0\% & 82.0\% & 86.7\% \\
    \bottomrule
  \end{tabular}
\end{table}

Two model inputs were corrected before the rates above were computed, and their
contributions separate cleanly. Measured object cross-sections replace first-pass
estimates for six objects---and, for the small screwdriver, the handle height is
separated from the \SI{15}{\milli\metre} width it is pinched across, which places
its grasp point at \SI{12}{\milli\metre} rather than \SI{7.5}{\milli\metre} and
is worth 5 iterative and 9 reflex trials---and the approach conditions use the
$\pm\SI{3}{\centi\metre}$ lateral jitter and elevated heights stated in the
hardware protocol in place of a $\pm\SI{150}{\milli\metre}$ reconstruction. With
the estimated sizes and the wide reconstruction the model gives 48.7\%, 65.3\%
and 77.3\%; correcting only the conditions gives 60.0\%, 84.7\% and 90.7\%;
correcting only the sizes gives 42.0\%, 65.3\% and 82.7\%. The conditions carry
iterative closure, because the $\pm\SI{150}{\milli\metre}$ starts require a
diagonal transit that drives the fingers through the support plane on the way in;
the sizes carry naive closure, by moving two more objects below the height at
which its terminal contact misses.
